\section{Discussion}
\label{sec:discussion}

The numbers above are read in this section rather than restated: what the
partition means, which sector of the model supplies the branch, how the
mechanism sits beside the published formulations, where every constant comes
from and what the model does not resolve.

\subsection{The observable: a modal partition of the scattered power}
\label{sec:disc_meaning}

The quantity at the centre of this article is $R_B$: the share of the incident
power that ends up in one particular outgoing mode, measured at a plane ten
thicknesses from the edge.  It is a partition of the scattered field rather than
a reflectance.  At the operating point the backward mode takes
\SI{1.93}{\percent} while the remaining propagating outgoing modes take
\SI{97.67}{\percent}; at the maximum the backward mode takes
\SI{94.29}{\percent} while the rest take \SI{5.63}{\percent}.  At the operating
point the total outgoing propagating power is $0.99980$ of the incident power
and at the maximum $0.99918$, so the plate is an almost perfect scatterer of the
guided mode in both cases and the interesting question is not how much is
reflected but into which branch it goes.

The distinction has a practical consequence.  A measurement of the total
reflectance of an edge would see essentially unity everywhere on the branch.  A
measurement that resolves the modal content -- by a wavenumber filter, or by a
time-of-flight separation of the branches -- would see the partition move from
almost nothing in the backward channel to more than half of the incident power.
The negative reflection discussed here is an observable that requires modal
resolution.

The word ``backward'' is used in an operational sense: a mode is backward when
the $x$ component of its thickness-integrated Poynting flux is negative, and
\eqref{eq:classification} confirms that the same mode is outgoing at the edge.
The dispersion branch has a region where $v_g$ and $v_E$ are negative
(\cref{fig:fig01}b), which is necessary for a backward mode to exist, and that
region is a property of the infinite plate.  What creates the conversion is the
edge, and what controls its strength is where the operating point sits on the
branch.

\subsection{The mechanism}
\label{sec:disc_mechanism}

The sweep along the branch (\cref{fig:fig04}) places the maximum of the
conversion at $k_xd=0.52719$, where the group velocity vanishes
($+0.45\,\si{\metre\per\second}$ against $-180$ and $+355\,\si{\metre\per\second}$
at the ends of the solved window) and where the backward wavenumber approaches
the incident one to within \SI{0.4}{\percent}.  This is a slow-mode effect and
it has a simple reading: at a zero-group-velocity fold the incident and backward
wavenumbers are nearly equal, so the two modes are nearly parallel and the edge
condition can be met with only a small admixture of the other outgoing modes;
the projection onto the backward mode then carries almost the whole partition.
Panel (c) of \cref{fig:fig04} makes this explicit: all coupling levels collapse
onto one curve in $k_{Bx}/k_{ix}$, with the conversion maximal as the ratio
approaches unity.  The same mechanism explains why the conversion falls steeply
on both flanks: on the fast side the wavenumbers separate, and on the far side of
the fold the branch has changed character.  A grid of $\Delta k_xd=0.05$ would miss this maximum, because the whole top of
the response is $10^{-3}$ wide; the two rasters of \cref{fig:fig04} locate it.

The reading has a second half.  At the fold the incident
mode itself is nearly cut off: its group velocity vanishes, so the flux it
carries per unit normalisation falls by a factor $128$ with respect to the
operating point.  The partition is maximal exactly where the transported power
is minimal.  A statement that ninety-five per cent of the energy is turned around
would therefore be true as a ratio and false as an energy budget, and the paper
makes only the first.

The reflection itself is present before the coupling is switched on.  The
controls of \cref{sec:results_mechanism} establish this directly: with $H=0$ the
same branch converts \SI{94.29}{\percent} of the incident power at the
maximum and \SI{1.93}{\percent} at the operating point, and when the
strain-gradient sector is removed the backward branch does not exist at all.
The strong negative reflection is therefore a property of the higher-order
guided branch and of the edge, not a resonance of the added coupling.

The coupling $H$ acts on the loss, and through the loss on the partition.  It
shifts the dispersion by at most $1.1\times10^{-5}$ in relative terms
(\cref{fig:fig01}a), it shifts the loss by up to a factor of sixty-two
(\cref{fig:fig06}c,d), and the partition follows with a weight that depends on
where the branch is sampled: \SI{-1.21}{\percent} at the operating point,
\SI{-10.17}{\percent} at $k_xd=0.55$ and \SI{-36.1}{\percent} at the maximum (\cref{tab:controls}).  The ordering --- dispersion unaffected,
partition affected only where the branch is slow, loss strongly affected
everywhere --- is what a carrier-level coupling acting through conduction loss
should produce.  The third of the partition lost at the maximum is mostly
loss acting on the backward mode between the edge and the reference plane:
referred to the edge plane, where both effects act at the same place, the same
comparison gives \SI{13.7}{\percent} rather than \SI{36.1}{\percent}, so
\SI{23.1}{} points are propagation loss and the rest is a genuine reduction of
the conversion.  Re-searching the extremum at each coupling level separates the
two effects further: the location moves from $0.52719$ to $0.5218$ and $0.5309$
at the two extremes of $H$, and the largest attainable share falls from $0.9468$
to $0.6244$ and $0.7013$, so the coupling both lowers the ceiling and displaces
it, asymmetrically in the sign of $H$.

The transverse wavenumber is the organising parameter at the operating point,
not the coupling.  \Cref{fig:fig07} shows that $k_y$ changes the backward
channel by a factor of $1.35$ across the solved box while $H$ moves it by
\SI{1.21}{\percent} branch-consistently (\SI{2.4}{\percent} at fixed frequency).
At the maximum, however, the coupling becomes the larger of the two
effects, moving the partition by \SI{36.1}{\percent} while the $k_y$ factor
stays at $1.35$.  The ordering is not a property of the coupling alone; it is a
property of the coupling and the operating point together.
Any statement about the added coupling must therefore be made at fixed $k_y$,
and any experimental attempt to see the coupling must control the transverse
wavenumber much better than it controls the doping or the mobility --- which is
consistent with the observation of \cref{fig:fig02} that changing the density by
a factor of three leaves the partition unchanged.

\subsection{Relation to the literature}
\label{sec:disc_literature}

The formulation belongs to the flexoelectric-semiconductor continuum developed
by Qu and co-workers \cite{qu2020mobile,qu2022thermal}, in which the strain
gradient polarises the medium, the polarisation is screened by the mobile
carriers, and the carrier redistribution feeds back on the field.  The wave
theory that we use is that of Yang et al.\ \cite{yang2018lamb,yang2020rayleigh}
for flexoelectric plates and half-spaces with strain-gradient elasticity, and the
carrier-physics conclusions of Wang et al.\ \cite{wang2025rayleigh} and Zhang et
al.\ \cite{zhang2025pscplate} --- that the semiconductor coupling is active only
within a window of carrier concentration and thickness --- are consistent with
our finding that $Q=d/L_D=0.075$ places the plate inside that window.

What is different here is the object of study.  The references above compute
dispersion relations and attenuation of modes in an unbounded or layered
medium, and the scattering calculations available for flexoelectric media are
plane-wave reflection coefficients at a free surface of a half-space
\cite{mondal2022scattering}.  This paper computes the \emph{scattering} of a
guided mode at a free edge of a plate, resolves the outgoing spectrum into its
branches, and reports the partition.  That is a different calculation with a different output: a
dispersion curve is an eigenvalue problem, whereas an edge partition requires
the matching of a continuum of outgoing modes and a classification of each of
them.

The closest primary work is the thermodynamically consistent theory of Qu, Pan,
Zhu and Deng \cite{qu2025thermodynamic}, who derive the constitutive equations
for a flexoelectric semiconductor from the free-energy imbalance with the
Coleman--Noll procedure, and who obtain the drift--diffusion current, the
flexoelectric and piezoelectric couplings and the higher-order stress in one
framework.  The primary article is closed access and no repository or preprint
copy is indexed, so its equations could not be inspected directly.  Its
free-energy
decomposition is, however, reproduced explicitly, with attribution, in the
open-access computational paper of Hu et al.\ \cite{hu2025computational}
(Eqs.~(4a)--(4e) there), and that reproduction is what we compare against.  In
that formulation the electric Gibbs free energy is additive,
\begin{equation}
\Psi_{\text{Gib}}=\Psi_{\text{Mech}}(\varepsilon,g)
+\Psi_{\text{Elec}}(e)+\Psi_{\text{Diele}}(\varepsilon,g,e)
+\Psi_{\text{Carrier}}(n,p),
\label{eq:qu_decomposition}
\end{equation}
with
$\Psi_{\text{Mech}}=\tfrac12C_{ijkl}\varepsilon_{ij}\varepsilon_{kl}
+\tfrac12A_{ijklmn}g_{ijk}g_{lmn}$,
$\Psi_{\text{Diele}}=-\eta_{ijk}e_i\varepsilon_{jk}
-\mu_{lijk}e_l g_{ijk}$ and
$\Psi_{\text{Carrier}}(n,p)$ a function of the two carrier densities alone.
Three consequences follow for the present paper.

First, the carrier energy in \eqref{eq:qu_decomposition} carries no strain
gradient.  Since the thermodynamic restriction fixes
$\partial\Psi_{\text{Gib}}/\partial n=\mu_n+q\psi$, the electrochemical potential
of that framework contains no strain-gradient contribution, and the strain
gradient reaches the carriers only \emph{indirectly}, through the flexoelectric
polarisation that enters Gauss's law.  The direct bilinear term
$\Delta n\,H_{jkl}\varepsilon_{jk,l}$ of \eqref{eq:added_coupling} is absent from
the postulate \eqref{eq:qu_decomposition}.  Within the equations we inspected it
is therefore an \emph{independent} addition, not a change of variables and not a
special case of a more general tensor; we do not claim that it is absent from
the full article, which we could not obtain, and we make no priority claim.

Second, the symbol $H$ is overloaded in the literature and we flag the collision
explicitly to avoid a misreading: in \cite{qu2025thermodynamic} and
\cite{hu2025computational}, $H_{ijklmn}$ is the sixth-order strain-gradient
elasticity tensor, which in this paper is written $A_{ijklmn}$ (\cref{eq:enthalpy}, \cref{eq:tau});
our $H_{jkl}$ is a third-order carrier--strain-gradient
coefficient with different units (J\,m).  The two are unrelated, and the
positivity bounds derived in \cref{sec:governing} apply only to the latter.

Third, the comparison locates the contribution.  If the part played by the
coupling were the headline, the contribution would be a constitutive hypothesis
of the kind \cite{qu2025thermodynamic} already systematises, and it is not the
headline: \cref{sec:results_mechanism} shows that the negative reflection exists
with $H=0$.  What this paper contributes is the \emph{scattering} problem --- the
partition of a guided mode into the outgoing branches of a free edge, with a
backward branch on which energy and phase travel in opposite directions ---
together with the demonstration that a carrier--strain-gradient enthalpy can be
added to a strain-gradient piezoelectric semiconductor without violating the
positivity bound, and a quantified account of its incremental effect.

The added term is therefore not put forward as new.  The formulations compared
above already contain a flexoelectric coupling of strain gradients to the
electric field, and the carrier density enters them through the free energy of
the electron gas; the term $\Delta n\,H_{jkl}\varepsilon_{jk,l}$ with its
reciprocal partners in $\sigma_{ij}$ and $\widetilde\mu$ is written here as the
minimal extension that closes the enthalpy, and the paper establishes as much
about what it does not do as about what it does: it moves the partition by
\SI{1.21}{\percent} at the operating point, \SI{9.85}{\percent} at
$k_xd=0.55$ and \SI{36.1}{\percent} at the maximum, all at $|H|=0.05H^{*}$, and
it is not the origin of the backward branch.  No priority claim is made for it.
The comparison is made against the equations that could be inspected, and the
full text of \cite{qu2025thermodynamic} was not among them; the question of
whether that formulation already contains a term of this structure therefore
remains open here rather than being settled in either direction.  What the
present numbers do establish is that the edge-scattering problem is well posed,
that its energy budget closes to $10^{-12}$, and that $\ell=0$ removes the
backward branch and with it the effect.

\subsection{Material provenance}
\label{sec:disc_material}

The provenance of the parameters is mixed and is stated explicitly in
\cref{tab:provenance}.  The classical constants are literature values; the
transport parameters are representative values for n-type ZnO at room
temperature linked by the Einstein relation; the strain-gradient length, the
flexoelectric coefficient and the coupling scale are \emph{synthetic}.  In particular, no measured value of $H_{jkl}$ exists, and the numbers used here
($|H|\le\SI{2.0e-23}{\joule\metre}$) are fixed by the positivity bound
\eqref{eq:schur}, which leaves the enthalpy strictly convex with a margin of a
factor of twenty; no specimen supplies them.

\begin{table}[t]
\centering
\caption{Provenance of every parameter, with the file in which it is recorded.}
\label{tab:provenance}
\small
\begin{tabular}{@{}p{0.2\textwidth}p{0.42\textwidth}p{0.3\textwidth}@{}}
\toprule
Parameter & Value and provenance & Source \\
\midrule
$c_{ij}$, $e_{ij}$, $\kappa_{ij}$, $\rho$ &
$c_{11}=210$, $c_{33}=221$, $c_{44}=43\,\si{\giga\pascal}$;
$e_{33}=1.32$, $e_{31}=-0.57$, $e_{15}=-0.48\,\si{\coulomb\per\square\metre}$;
$\kappa_{33}=8.85\times10^{-11}\,\si{\farad\per\metre}$; literature values &
\cite{zhang2025pscplate}, Table~1;
\texttt{a0/materials.json} \\
$\mu_e$, $D$ & $0.0175\,\si{\square\metre\per\volt\per\second}$,
$D=\mu_e k_BT/q=4.55\times10^{-4}\,\si{\square\metre\per\second}$; representative &
Einstein relation, $T=\SI{300}{\kelvin}$ \\
$n_0$ & $1\text{--}3\times10^{18}\,\si{\per\cubic\metre}$; modelling scenario &
\texttt{fig02} sweep \\
$d$ & \SI{200}{\nano\metre} reference, 120--\SI{320}{\nano\metre} sweep;
modelling scenario & \texttt{fig03} sweep \\
$\ell$ & \SI{50}{\nano\metre}; \textbf{synthetic} & this work \\
$f_{3zz}$ & \SI{1}{\nano\coulomb\per\metre}; \textbf{synthetic}, contributes
\SI{-0.01}{\percent} to $R_B$ & this work \\
$H_{jkl}$ & $H_{333}$ only, $|H|\le2.0\times10^{-23}\,\si{\joule\metre}
=0.022\,H_{\text{bound}}$; \textbf{synthetic} & this work \\
\bottomrule
\end{tabular}
\end{table}

\subsection{Limitations}
\label{sec:disc_limitations}

Five limitations bound the conclusions, and each is stated where it enters the
model.

\begin{enumerate}
\item \emph{Geometry.}  The formulation is two-and-a-half-dimensional: a plate
infinite in $y$, semi-infinite in $x$, with three displacement components on the
$(x,z)$ cross-section.  It has one edge and not two, and it does not resolve the
three-dimensional corner where the edge meets the free surfaces.  A finite
plate, a second edge, and the interference between them are outside the model.
\item \emph{The corner.}  Corner effects were assessed rather than assumed, in
the only way this formulation allows: the strain-gradient energy density on the
edge plane was computed from the coherent sum of the incident mode and all
outgoing modes, and its concentration in the outer band $|z|>0.4\,d/2$ was
followed through the degree refinement.  At the operating point the band carries
$0.93$--$1.00$~\% of that energy, at the maximum $1.22$--$1.30$~\% and at
$k_xd=0.55$ $1.18$--$1.26$~\%, its peak-to-mean ratio inside the band is
$0.098$--$0.102$, $0.126$--$0.130$ and $0.122$--$0.126$ (below the mean), and its
share \emph{decreases} by \SI{6.8}{}, \SI{6.6}{} and \SI{6.6}{}~\% between degree
18 and degree 30 instead of growing.  A corner singularity would show the
opposite trend, so the reported partitions (degree spread $7.2\times10^{-7}$ and
$1.5\times10^{-7}$) are not supported by a corner energy reservoir.  That is
evidence about the energy and not a proof that the corner line force is
dynamically negligible: a small band fraction bounds the energy near the corner
in the computed field, but it does not bound the work a singular corner traction
could do on the admissible variations, because the polynomial basis cannot
represent that singularity in the first place.  The global partitions are
insensitive to the corner at the level shown; the pointwise corner field and the
line force associated with it remain unresolved, and settling them would need a
corner-eigenvalue analysis of the higher-order edge problem or corner-enriched
trial functions, neither of which is attempted here.
\item \emph{The carrier model and the electrical condition.}  The carrier
physics is linearised about a uniform doping and zero bias, so the nonlinear
drift--diffusion regime, in which the perturbation is not small compared with
$n_0$, is outside the model.  The open-circuit condition required a gauge to
remove the null space of the electrostatic operator; the gauge changes no
physical field, but it also means that the open-circuit results are a solvable
version of the same physics rather than an independent check of it.
\item \emph{The parameters.}  The classical constants are those of a bulk
crystal and the transport parameters are representative values for n-type ZnO at
room temperature; the strain-gradient length, the flexoelectric coefficient and
the coupling scale are synthetic, and no measured value of $H_{jkl}$ exists.
The reported percentages should therefore be read as the response of this model
and not as predictions for a particular film, and the flexoelectric sector, which
is present but subdominant (\cref{sec:results_scope}), is not a calibrated
mechanism.
\item \emph{The constitutive attribution.}  The relation of the added coupling
to the published flexoelectric--semiconductor formulations remains open, as
discussed in \cref{sec:disc_literature}; no priority is claimed for it.  What the
model does establish is that a bilinear carrier--strain-gradient enthalpy can be
added inside the positivity bound, that its reciprocal partners follow from the
same potential, and that its incremental effect on the partition is measurable,
of one sign, and small compared with the effect of the branch itself.
\end{enumerate}
